\documentclass{beamer}
\usetheme{Madrid}
\usecolortheme{whale}

\usepackage[utf8]{inputenc}
\usepackage[spanish]{babel}
\usepackage{amsmath}
\usepackage{amsfonts}
\usepackage{amssymb}

% --- Título y Datos ---
\title{Álgebra de Heyting}
\subtitle{Retícula brouweriana}
\author{Karina G. Buendía}
\date{abril de 2026}

\begin{document}

% Slide 1: Portada
\frame{\titlepage}

% Slide 2: Introducción - Retículas
\begin{frame}{Retículas}
    Una retítucla $L$ es un conjunto parcialmente ordenado que tienen dos productos binarios. Si una retícula $L$ tiene
    elementos $0$ y $1$ se puede definir ecuacionalmente como un conjunto con dos elementos distinguidos y dos operaciones
    binarias $\land$ y $\lor$; estas operaciones son tanto asociativas como conmutativas \dots
\end{frame}

% Slide 3: Introducción - Álgebra de Boole
\begin{frame}
    \dots que satisfacen las siguientes identidades.

    \begin{enumerate}
        \item $a \land a = a$, $x \lor a =a$,
        \item $1 \land a = a$, $0 \lor a =a$,
        \item $a \land (b \lor a) = a = (a \land b) \lor a$. 
    \end{enumerate}
\vspace{0.5cm}
    Dadas estas ecuaciones, el orden parcial puede ser recuperado porque $a \leq b$ si y solo si $a \land b = a$ (o bien,
    de manera equivalente $a \lor b = b$).

\vspace{0.5cm}
    El conjunto parcialmente ordenado de todos los subconjuntos de un conjunto dado es siempre un álgebra de booleana.
\end{frame}

\begin{frame}{Álgebra de Heyting}
    Un álgebra de Heyting $H$ (también llamado retícula brouweriana) es una retícula con $0$ y $1$. Dicha retícula tiene,
    para cualquier par de elementos, $x$ y $y$ un exponencial $y^x$. Este exponencial es usualmente denotado como
    $x \rightarrow y$, que es la mínima cota superior de todos quellos elementos $z \land x \leq y$. Es decir, $z \leq
    (x\rightarrow y) \iff z \land x \leq y$.
    \begin{enumerate}
        \item Al realizar la intersección con $x$ el resultado no pasa de $y$.
        \item Cada elemento $z$ que esté por debajo de $y$ (con el ínfimo) debe ser menor o igual al elemento implicación.
    \end{enumerate}
\end{frame}

\end{document}